\documentclass{article}

% if you need to pass options to natbib, use, e.g.:
%     \PassOptionsToPackage{numbers, compress}{natbib}
% before loading neurips_2026

% The authors should use one of these tracks.
% Before accepting by the NeurIPS conference, select one of the options below.
% 0. "default" for submission
\usepackage{neurips_2026}
% the "default" option is equal to the "main" option, which is used for the Main Track with double-blind reviewing.
% 1. "main" option is used for the Main Track
%  \usepackage[main]{neurips_2026}
% 2. "position" option is used for the Position Paper Track
%  \usepackage[position]{neurips_2026}
% 3. "eandd" option is used for the Evaluations & Datasets Track
 % \usepackage[eandd]{neurips_2026}
% 4. "creativeai" option is used for the Creative AI Track
%  \usepackage[creativeai]{neurips_2026}
% 5. "sglblindworkshop" option is used for the Workshop with single-blind reviewing
 % \usepackage[sglblindworkshop]{neurips_2026}
% 6. "dblblindworkshop" option is used for the Workshop with double-blind reviewing
%  \usepackage[dblblindworkshop]{neurips_2026}
% 7. "education" option is used for the Education Track
%  \usepackage[education]{neurips_2026}


% After being accepted, the authors should add "final" behind the track to compile a camera-ready version.
% 1. Main Track
 % \usepackage[main, final]{neurips_2026}
% 2. Position Paper Track
%  \usepackage[position, final]{neurips_2026}
% 3. Evaluations & Datasets Track
 % \usepackage[eandd, final]{neurips_2026}
% 4. Creative AI Track
%  \usepackage[creativeai, final]{neurips_2026}
% 5. Workshop with single-blind reviewing
%  \usepackage[sglblindworkshop, final]{neurips_2026}
% 6. Workshop with double-blind reviewing
%  \usepackage[dblblindworkshop, final]{neurips_2026}
% Note. For the workshop paper template, both \title{} and \workshoptitle{} are required, with the former indicating the paper title shown in the title and the latter indicating the workshop title displayed in the footnote.
% For workshops (5., 6.), the authors should add the name of the workshop, "\workshoptitle" command is used to set the workshop title.
% \workshoptitle{WORKSHOP TITLE}
% 7. Education Track
% \usepackage[education, final]{neurips_2026}

% "preprint" option is used for arXiv or other preprint submissions
 % \usepackage[preprint]{neurips_2026}

% to avoid loading the natbib package, add option nonatbib:
%    \usepackage[nonatbib]{neurips_2026}

\usepackage[utf8]{inputenc} % allow utf-8 input
\usepackage[T1]{fontenc}    % use 8-bit T1 fonts
\usepackage{hyperref}       % hyperlinks
\usepackage{url}            % simple URL typesetting
\usepackage{booktabs}       % professional-quality tables
\usepackage{amsfonts}       % blackboard math symbols
\usepackage{nicefrac}       % compact symbols for 1/2, etc.
\usepackage{microtype}      % microtypography
\usepackage{xcolor}         % colors
\usepackage{float}
\usepackage{amsmath}
\usepackage{graphicx}

\title{Feature Learning in Supervised AI Models}


\author{
  Zain Farhan \\
  Lahore University of Management Sciences (LUMS) \\
  \texttt{Department of Computer Sciences} \\
}


\begin{document}


\maketitle


\begin{abstract}
  In this research paper we will learn and investigate how various supervised learning models learn features. We will do this
  through four tasks. In task 1 we will see how strong inductive bias is in models and how their architecture can effect how they
  learn features (i.e. whether they learn shape or texture or some other attribute). In task 2 we will observe how this effects Domain
  Generalization. The idea being that we want the model to learn identifying features about the objects (like their shape) instead of Domain
  Specific features (like if the dataset had grass behind every cat and the model associated the grass with the cat). If the model can do this, it should
  perform well even on Domain Shifts. We extend this idea in task 3 with Domain Generalization and completely unknown target domains. Finally, in task 4
  we tackle the open set problem, to see if we can identify casual features even on unknown classes, and separate them as they come.
\end{abstract}

\section{Introduction}
\label{sec:intro}

We will do various experiments in this research paper to study how different AI models learn features, and the consequence of their architectures.
This is done in four parts.

First, we will see the study the influence of Inductive Biases and Feature Representations, our goal here is to see the influence of model architecture.
By comparing the metrics of four well known and standard models (ResNet-50, ViT-B, CLIP). We will first perform a baseline metric on all of them, and then
perform 5 different experiments. Colour Bias: We perform two different types, grey scale, and hue shift, and see how well the models perform.
Shape vs Texture (cue-conflict): We apply the texture of one object to another, to see if models are using texture or shape for classification and which is more important.
Translation: We shift the images, left, right, up, or down and see if the model can still classify. Lastly, we shuffle the patches of the images
to see how the global structure effects the models.

Our points of interest are, what cue-conflict and colour reveal about reliance on shape, texture, and colour, and what coverage can tell us about shape bias.
How the models respond when spatial structure is altered and what that tells us about the importance of locality, global organisation, and positional sensitivity.
To what extent prediction-level changes agree with changes in feature representations, and which observed difference can reasonably be attributed
to architecture as opposed to pretraining data, supervision, augmentation or model capacity.

\section{Experimental Setup/Methods}
\label{sec:setup}

\subsection{Inductive Biases and Feature Representations}

\textbf{Models:}
\begin{itemize}
  \item \texttt{torchvision ResNet-50 with ResNet50\_Weights.IMAGENET1K\_V2;}
  \item \texttt{torchvision ViT-B/16 with ViT\_B\_16\_Weights.IMAGENET1K\_V1;}
  \item \texttt{OpenCLIP ViT-B-32 with pretrained='openai'.}
\end{itemize}

\textbf{Dataset:} STL-10

\textbf{Configuration:}
\begin{itemize}
  \item Seed: 6304
  \item Optimizer: AdamW
  \item Learning Rate: $10^{-3}$
  \item Weight Decay: $10^{-4}$
\end{itemize}

For the Shape vs Texture we used AdaIN implementation, which follows Huang \&
Belongie (2017). the encoder/decoder architecture and the pretrained weights come from the public reimplementation
naoto0804/pytorch-AdaIN, weights mirrored at BhupatiNadar/adain-style-transfer-models
on HuggingFace.



\section{Results}
\label{sec:results}

\subsection{Inductive Biases and Feature Representations}

We first take some basic measurements as a baseline, we freeze the three classifier heads and measure top-1 accuracy,
macro-F1, and mean maximum confidence. We also do this for zero shot CLIP, we measure confidence from the softmax over scaled
class similarities.

\begin{table}[H]
  \caption{Clean Baseline - Shows all models performing very well on the dataset.}
  \label{tab:clean-baseline}
  \centering
  \begin{tabular}{lccc}
    \toprule
    System & Top-1 (\%) & Macro-F1 (\%) & Mean max conf. \\
    \midrule
    ResNet-50 + linear head   & 96.4 & 96.39 & 0.91 \\
    ViT-B/16 + linear head    & 97.2 & 97.21 & 0.97 \\
    CLIP ViT-B-32 + linear head & 96.6 & 96.60 & 0.27 \\
    CLIP ViT-B-32 zero-shot   & 97.0 & 97.01 & 0.94 \\
    \bottomrule
  \end{tabular}
\end{table}

All the models perform well as a baseline, this maybe a limitation of the STL-10 dataset, as its distinct classes would be easier to classify.
This can be a limitation as differences in experiments would more visibly show a decrease rather than an increase.
It should also be noted that the CLIP Max Mean is low because it uses Cosine similarity in training, and the L-2 norm forces the magnitude to be 1.
So when attached with a linear head and no compensation, the value comes off as low. It does not mean it is lower in confidence than other models,
just that rows can't be compared the magnitude lower, not because it is less confident than others

\subsubsection{Color Bias}

We did two experiments with colour, the first we changed the images to greyscale measured the same metrics and got the following results:

\begin{table}[H]
  \caption{Greyscale: Minimal change in Top-1 Accuracy}
  \label{tab:grayscale}
  \centering
  \begin{tabular}{lccccc}
    \toprule
    System & Top-1 & Macro-F1 & $\Delta$ Top-1 & Consistency & Conf. \\
    \midrule
    ResNet-50 + linear head     & 94.6 & 94.61 & $-1.8$ & 95.8 & 0.91 \\
    ViT-B/16 + linear head      & 95.4 & 95.41 & $-1.8$ & 98.0 & 0.96 \\
    CLIP ViT-B-32 + linear head & 93.8 & 93.78 & $-2.8$ & 93.0 & 0.25 \\
    CLIP ViT-B-32 zero-shot     & 93.0 & 92.97 & $-4.0$ & 94.0 & 0.93 \\
    \bottomrule
  \end{tabular}
\end{table}

We can see that making the image greyscale results in a slight decrease in accuracy throughout. Given that we have removed information from the image
this is to be expected. A more interesting observation is that there was no significant difference between Convolution and Vision Transformers, 
though consistency is slightly higher for the Vision Transformer. CLIP had the most sever impact. Both of the readings are lower than either
ResNet or ViT-B, and further zero-shot is the lowest. One plausible explanation for this is that the text embeddings had an influence on how
accurately important colour is during pre training. For example if a prompt were to be "A picture of a Brown Dog." that would influenced both the visual
encoder and the text encoder. It would also explain why zero-shot performed even worse. This though is not conclusive, and would be hard to confirm.

The second experiment was with a Hue Shift of $40^{}$ deg. A limitation of this is that a hue shift will only have a good effect on well
saturated images. On less saturated images the change is very minimal, making the result less significant.

\begin{table}[H]
  \caption{Hue Rotation of 40: Minimal change in accuracy}
  \label{tab:hue40}
  \centering
  \begin{tabular}{lccccc}
    \toprule
    System & Top-1 & Macro-F1 & $\Delta$ Top-1 & Consistency & Conf. \\
    \midrule
    ResNet-50 + linear head     & 94.4 & 94.42 & $-2.0$ & 95.2 & 0.90 \\
    ViT-B/16 + linear head      & 95.8 & 95.79 & $-1.4$ & 98.2 & 0.96 \\
    CLIP ViT-B-32 + linear head & 94.8 & 94.78 & $-1.8$ & 95.8 & 0.26 \\
    CLIP ViT-B-32 zero-shot     & 95.2 & 95.18 & $-1.8$ & 95.6 & 0.93 \\
    \bottomrule
  \end{tabular}
\end{table}

Again we see a minimal change in accuracy, this time they are much closer to each other. No claims can be made as the difference is small enough that it can
be considered noise. Again though, we will find that ViT is just slightly more consistent than other models. It is also of note that unlike greyscale, CLIP is
much more in line with other models. One explanation for this is that in greyscale, we totally lost colour information, however with a rotation, the relationship
between the colour of the objects is still the same, it has just been rotated. Hence, the loss is less with CLIP now.

\subsubsection{Shape vs. Texture}

For this comparison alpha strength of 1 was required, as lower values did not transfer texture well enough. We get the following results:

\begin{table}[H]
  \caption{Shape vs Texture: Convoutional models are not shape biased.}
  \label{tab:shape-texture-both}
  \centering
  \begin{tabular}{lcccrrrr}
    \toprule
    & \multicolumn{5}{c}{Balanced ($n=220$)} & \multicolumn{2}{c}{All accepted ($n=286$)} \\
    \cmidrule(lr){2-6} \cmidrule(lr){7-8}
    System & $N_{\text{sh}}$ & $N_{\text{tex}}$ & $N_{\text{oth}}$ & Shape bias & Coverage & Shape bias & Coverage \\
    \midrule
    ResNet-50 + linear head     &  54 & 52 & 114 & 50.94 & 48.18 & 57.78 & 47.20 \\
    ViT-B/16 + linear head      &  99 & 14 & 107 & 87.61 & 51.36 & 89.94 & 55.59 \\
    CLIP ViT-B-32 + linear head & 116 & 12 &  92 & 90.62 & 58.18 & 91.72 & 59.09 \\
    CLIP ViT-B-32 zero-shot     & 117 &  7 &  96 & 94.35 & 56.36 & 95.81 & 58.39 \\
    \bottomrule
  \end{tabular}
\end{table}

Out of a total of 400 candidate images only 286 were of good enough quality, which were further reduced to 220 to make each class balanced.
As expected, the Vision Transformers have a considerable shape bias, while the convolutional model does not. Though one point of note is that
ResNet still managed to score a 50.94 on the Shape Bias. This is considerably higher than the original Geirhos et al. (2019) paper. The reason for this
can be mostly likely presumed to be the biggest limitation of this experiment.

That is that the images were filtered out manually. This introduced significant human bias in the process making the data unreliable.
The human eye is considered to be heavily biased towards shape, I personally noticed as well that during the filtering it was much easier for me to tell when the shape data was destroyed
in comparison to when texture data was to sufficiently applied. It is reasonable to then conclude the data had a bias towards higher quality shape images than texture images.

Another notable observation is that for ViT and ResNet both, there most quantitative category is neither shape nor texture, and CLIP barely has more shape classified images.
Making the coverage very low, and limiting our ability to draw further conclusions from this data. This can again be contributed to the low quality of the data. The AdaIN
implementation resulted in images with quite high distortion with our particular dataset, hampering the ability of the models to classify.

\subsubsection{Patch Structure}

We divided the images into 4x4 patches, and shuffled the patches. We feed the shuffled images back in the model to measure the metrics:

\begin{table}[H]
  \caption{Patch Shuffling: Transformers suffer the most}
  \label{tab:patch-shuffle}
  \centering
  \begin{tabular}{lccccc}
    \toprule
    System & Top-1 (\%) & Macro-F1 (\%) & $\Delta$ Top-1 & Consistency (\%) & Mean max conf. \\
    \midrule
    ResNet-50 + linear head     & 89.0 & 89.06 & $-7.4$  & 90.4 & 0.87 \\
    ViT-B/16 + linear head      & 90.4 & 90.52 & $-6.8$  & 92.6 & 0.89 \\
    CLIP ViT-B-32 + linear head & 84.4 & 84.18 & $-12.2$ & 84.6 & 0.20 \\
    CLIP ViT-B-32 zero-shot     & 84.4 & 84.20 & $-12.6$ & 84.2 & 0.82 \\
    \bottomrule
  \end{tabular}
\end{table}

When we shuffle the images we destroy the global structure of the image, i.e. the shape. For Vision Transformers this is quite bad as they use the self attention mechanism and look at a global picture.
Understandably, they suffer the most in accuracy. Convolutional networks however only look at the local area, making them much more resilient to this, further
texture is preserved more, as each patch can act like a texture patch and we saw in table\ref{tab:shape-texture-both}, they are much more biased towards texture.

The ViT-B and ResNet model are close in accuracy. This is a bit unexpected because we discussed ResNet should perform noticeably better here. The likely culprit is that we used zero padding,
max pooling, stride 2, and we averaged the final 7x7 map which all contribute to either a lack of local information or hampers its readability for the models.

One final thing of note is that CLIP performs noticeably worse than ViT, even though they share the visual encoder (as a Vision Transformer), the fact that both zero-shot and linear head land at the same number
implies that its not the involvement of the text encoder in the prediction process. This suggest that contrastive learning (with the text encoder) during pre training or some other stage had some factor that made
them more dependent factor.\footnote{We also saw more shape bias for them in table\ref{tab:shape-texture-both}, but it was not enough of a difference to warrant any conclusion}

\subsubsection{Representation Analysis}

We have discussed most experiments now, but those mainly measured when the predictions changed and hence also changed the accuracy. But how the model sees and interprets the feature can change even if the
prediction didn't. We measure the Cosine Similarity and t-SNE maps where relevant.

\begin{table}[H]
  \caption{Cosine similarities - Expected Results}
  \label{tab:representation-stability}
  \centering
  \small
  \begin{tabular}{lcccccc}
    \toprule
    & \multicolumn{2}{c}{ResNet-50} & \multicolumn{2}{c}{ViT-B/16} & \multicolumn{2}{c}{CLIP ViT-B-32} \\
    \cmidrule(lr){2-3} \cmidrule(lr){4-5} \cmidrule(lr){6-7}
    Intervention & $I_T$ & Norm. & $I_T$ & Norm. & $I_T$ & Norm. \\
    \midrule
    Grayscale                  & 0.789 & 0.719 & 0.737 & 0.711 & 0.900 & 0.743 \\
    Translation ($32$\,px)     & 0.934 & 0.912 & 0.943 & 0.938 & 0.970 & 0.923 \\
    Patch shuffle              & 0.718 & 0.626 & 0.649 & 0.615 & 0.728 & 0.299 \\
    Cue conflict vs.\ content  & 0.250 & 0.004 & 0.223 & 0.146 & 0.651 & 0.101 \\
    Cue conflict vs.\ style    & 0.216 & $-0.041$ & 0.072 & $-0.020$ & 0.580 & $-0.083$ \\
    \midrule
    Random-pair floor          & 0.247 & 0.000 & 0.090 & 0.000 & 0.612 & 0.000 \\
    \bottomrule
  \end{tabular}
\end{table}

In terms of grayscale and translation we can see the Norm is rather high. This is consistent with our earlier hypothesis that
models don't look at colour for classification. And since in translation most features are preserved (shape, texture, colour) and the image in only moved, they are very close to the original
in representation.

In Patch shuffling the values a bit lower (around 0.6), which is to be expected as the global structure is destroyed. What is more interesting is that CLIP and ViT show a large difference despite
them both having a Vision Transformer. One plausible explanation is that since CLIP's tokens are twice as large as ViT, they are more dependent on global
structure, and the representation drifts more. The larger token size could also mean that it covers the boundaries more between patches.

The t-SNE diagrams (Figures\ref{fig:tsne-clip} \ref{fig:tsne-resnet} \ref{fig:tsne-vit} in the appendix) agree with these results.

\subsection{Unsupervised Domain Adaptation}

\subsubsection{Source Only ERM}

We first train both the model and the classifier hear on the source domains using cross entropy, which results in these measurements.

\begin{table}[H]
  \caption{Source only ERM - the main domain gap identified}
  \label{tab:source-only}
  \centering
  \begin{tabular}{lcc}
    \toprule
    Domain & Accuracy (\%) & Macro-F1 (\%) \\
    \midrule
    Photo (source, val)         & 97.30 & 96.85 \\
    Art Painting (source, val)  & 91.71 & 91.69 \\
    Cartoon (source, val)       & 95.74 & 96.00 \\
    \midrule
    Mean source                 & 94.91 & 94.85 \\
    Worst source                & 91.71 & 91.69 \\
    \midrule
    Sketch (target)             & 61.59 & 63.09 \\
    \bottomrule
  \end{tabular}
\end{table}

As we can see in source only statistics, the accuracy on our source domains is very good, but when ran on a unknown domain (for us that is sketch) we notice a massive accuracy drop.
compared with the mean it is a -33 \% drop, this is the main figure we try to lower in these experiments.

\begin{table}[H]
  \caption{Target Domain by class - high amounts of variation.}
  \label{tab:source-only-per-class}
  \centering
  \begin{tabular}{lc}
    \toprule
    Class & Sketch accuracy (\%) \\
    \midrule
    Dog      & 21.11 \\
    Elephant & 94.73 \\
    Giraffe  & 77.95 \\
    Guitar   & 70.07 \\
    Horse    & 46.32 \\
    House    & 83.75 \\
    Person   & 61.25 \\
    \bottomrule
  \end{tabular}
\end{table}

when we see the by class accuracy on the target domain, dog performs worst while elephant performs best.
Dog is the largest class in the sketch data, we found dog's 333 most common errors go to elephant, 167 to giraffe. It is probable it is not confusing
animals and objects at least, it is confused within animal classes.

\subsubsection{DAN - MMD Alignment}

\begin{table}[H]
  \caption{DA with MMD - Slight improvement over Source Only}
  \label{tab:dan}
  \centering
  \begin{tabular}{lcc}
    \toprule
    & Source-only & DAN ($\lambda_{\mathrm{MMD}}=1$) \\
    \midrule
    Photo val acc (\%)          & 97.30 & 95.80 \\
    Art Painting val acc (\%)   & 91.71 & 93.41 \\
    Cartoon val acc (\%)        & 95.74 & 95.95 \\
    \midrule
    Mean source acc (\%)        & 94.91 & 95.05 \\
    Mean source macro-F1 (\%)   & 94.85 & 94.93 \\
    \midrule
    Sketch acc (\%)             & 61.59 & 64.44 \\
    Sketch macro-F1 (\%)        & 63.09 & 60.55 \\
    $\Delta$ Sketch acc         & ---   & $+2.85$ \\
    Domain separability (\%)    & 100.00 & 81.87 \\
    \bottomrule
  \end{tabular}
\end{table}

If we observe the general by class results for MMD, we see a slight but still meaningful improvement over the sketch accuracy.
Though the sketch F1 - macro decreases by a bit.

\begin{table}[H]
  \caption{MMD by Class - Good by class improvement}
  \label{tab:dan-per-class}
  \centering
  \begin{tabular}{lccc}
    \toprule
    Class & Source-only (\%) & DAN (\%) & $\Delta$ \\
    \midrule
    Dog      & 21.11 &  47.93 & $+26.82$ \\
    Elephant & 94.73 &  75.54 & $-19.19$ \\
    Giraffe  & 77.95 &  61.62 & $-16.33$ \\
    Guitar   & 70.07 &  74.18 & $+4.11$ \\
    Horse    & 46.32 &  57.60 & $+11.28$ \\
    House    & 83.75 & 100.00 & $+16.25$ \\
    Person   & 61.25 &  86.25 & $+25.00$ \\
    \bottomrule
  \end{tabular}
\end{table}

We see all classes improve in accuracy, except for giraffe and elephant, but given that in source only most of the errors
in the dog class went to giraffe and elephant, it is likely they were over represented in the original model in the first place.

\subsubsection{DANN - Adversarial Alignment}

We now use adversarial learning to see how well it performs. And get the following results:

\begin{table}[H]
  \caption{DANN - The performances worsened}
  \label{tab:dann}
  \centering
  \begin{tabular}{lcc}
    \toprule
    & Source-only & DANN \\
    \midrule
    Photo val acc (\%)          & 97.30 & 95.80 \\
    Art Painting val acc (\%)   & 91.71 & 92.93 \\
    Cartoon val acc (\%)        & 95.74 & 93.60 \\
    \midrule
    Mean source acc (\%)        & 94.91 & 94.11 \\
    Mean source macro-F1 (\%)   & 94.85 & 94.03 \\
    \midrule
    Sketch acc (\%)             & 61.59 & 55.97 \\
    Sketch macro-F1 (\%)        & 63.09 & 58.23 \\
    $\Delta$ Sketch acc         & ---   & $-5.62$ \\
    Domain separability (\%)    & 100.00 & 99.31 \\
    \bottomrule
  \end{tabular}
\end{table}

Interestingly not only is the performance on the sketch pictures worse than source, but the the predictions on source only is also worse.
At first this might make it look like DANN didn't work, but after further investigation we find this:

\begin{table}[t]
  \caption{Tracking Per Epoch - Loss was decreasing while Domain accuracy was going to 0.5}
  \label{tab:dann-training}
  \centering
  \begin{tabular}{ccccc}
    \toprule
    Epoch & Cls.\ loss & Domain loss & Domain acc. & Mean source F1 (\%) \\
    \midrule
    1 & 0.533 & 0.384 & 0.817 & 91.37 \\
    2 & 0.234 & 0.597 & 0.673 & 89.19 \\
    3 & 0.159 & 0.669 & 0.597 & \textbf{94.03} \\
    4 & 0.261 & 0.742 & 0.542 & 84.62 \\
    5 & 0.176 & 0.686 & 0.551 & 92.80 \\
    6 & 0.084 & 0.681 & 0.541 & 92.37 \\
    7 & 0.070 & 0.694 & 0.502 & 93.92 \\
    8 & 0.081 & 0.694 & 0.501 & 93.90 \\
    \bottomrule
  \end{tabular}
  \\[2pt]
  \footnotesize Selected epoch in bold. Chance domain accuracy is 0.5 and $\ln 2 = 0.693$.
\end{table}

The Loss was indeed decreasing, and the Domain Accuracy was going to 0.5., what this means is that even if we fool the discriminator.
that does not necessarily mean an increase in performance. Perhaps the backbone managed to fool the discriminator before the discriminator
got good enough to warrant an improvement. An increase in learning later or even just running more Epochs could improve metrics to match MMD or beyond. 

\subsubsection{CDAN - Class-Conditional Adversarial Alignment}

To improve the scores from before we use the values of the prediction as well, we do a cross product with the latent and feed that to the discriminator.

\begin{table}[t]
  \caption{CDAN: Significant Improvement from before}
  \label{tab:cdan}
  \centering
  \begin{tabular}{lcc}
    \toprule
    & Source-only & CDAN \\
    \midrule
    Photo val acc (\%)          & 97.30 & 97.30 \\
    Art Painting val acc (\%)   & 91.71 & 91.46 \\
    Cartoon val acc (\%)        & 95.74 & 95.52 \\
    \midrule
    Mean source acc (\%)        & 94.91 & 94.76 \\
    Mean source macro-F1 (\%)   & 94.85 & 94.66 \\
    \midrule
    Sketch acc (\%)             & 61.59 & 74.65 \\
    Sketch macro-F1 (\%)        & 63.09 & 73.85 \\
    $\Delta$ Sketch acc         & ---   & $+13.06$ \\
    Domain separability (\%)    & 100.00 & 99.86 \\
    \bottomrule
  \end{tabular}
\end{table}

The performance on the source only metric stays the same, but the accuracy on sketch images received a significant improvement.
Meaning that even soft information about the final prediction was very helpful to the discriminator. Improving its performance which
improved the backbone to more aggressively compete with it. However, it was not lower domain separability that did this, as domain separability
essentially remained the same. It can be speculated that domain information is important as that helps the model to classify them.

\subsubsection{Common Evaluation and Alignment Diagnostic}

We compare the results overall:

\begin{table}[H]
  \caption{Diagnostic - Separability might not be the deciding factor}
  \label{tab:t2-diagnostic}
  \centering
  \footnotesize
  \begin{tabular}{lccccccccc}
    \toprule
    & \multicolumn{3}{c}{Source val.\ acc.\ (\%)} & \multicolumn{2}{c}{Mean source (\%)} & \multicolumn{2}{c}{Sketch (\%)} & & \\
    \cmidrule(lr){2-4} \cmidrule(lr){5-6} \cmidrule(lr){7-8}
    Method & Photo & Art & Cartoon & Acc. & F1 & Acc. & F1 & $\Delta$ Acc. & Sep.\ (\%) \\
    \midrule
    Source-only & 97.30 & 91.71 & 95.74 & 94.91 & 94.85 & 61.59 & 63.09 & ---       & 100.00 \\
    DAN         & 95.80 & 93.41 & 95.95 & 95.05 & 94.93 & 64.44 & 60.55 & $+2.85$   &  81.87 \\
    DANN        & 95.80 & 92.93 & 93.60 & 94.11 & 94.03 & 55.97 & 58.23 & $-5.62$   &  99.31 \\
    CDAN        & 97.30 & 91.46 & 95.52 & 94.76 & 94.66 & 74.65 & 73.85 & $+13.06$  &  99.86 \\
    \bottomrule
  \end{tabular}
  \\[2pt]
  \footnotesize Sep.\ = domain separability: held-out accuracy of a logistic-regression probe
  ($C=1$, balanced, 70/30 split, seed 6304) trained to tell source from target from the frozen
  features. Chance is 50\%. Target labels were used only at this stage.
\end{table}

We would assume that bringing the domains closer and making them indistinguishable would be key, as the overlapping elements between the target domain and
source domains would come into play. However aside from MMD which tries to directly to do this with the Loss function, a freshly trained logistical regressor can
still separate the classes. And our best performance from CDAN barely shows any change, implying it might not be the most key factor.

\begin{table}[H]
  \caption{Per Class Comparison: Four classes share a trend}
  \label{tab:t2-per-class}
  \centering
  \small
  \begin{tabular}{lccccccc}
    \toprule
    & \multicolumn{4}{c}{Sketch accuracy (\%)} & \multicolumn{3}{c}{$\Delta$ vs Source-only} \\
    \cmidrule(lr){2-5} \cmidrule(lr){6-8}
    Class & Source-only & DAN & DANN & CDAN & DAN & DANN & CDAN \\
    \midrule
    Dog      & 21.11 &  47.93 &  42.10 &  64.25 & $+26.82$ & $+20.99$ & $+43.14$ \\
    Elephant & 94.73 &  75.54 &  35.00 &  80.54 & $-19.19$ & $-59.73$ & $-14.19$ \\
    Giraffe  & 77.95 &  61.62 &  51.53 &  62.95 & $-16.33$ & $-26.42$ & $-15.00$ \\
    Guitar   & 70.07 &  74.18 &  98.36 &  89.31 & $+4.11$  & $+28.29$ & $+19.24$ \\
    Horse    & 46.32 &  57.60 &  48.04 &  71.69 & $+11.28$ & $+1.72$  & $+25.37$ \\
    House    & 83.75 & 100.00 & 100.00 & 100.00 & $+16.25$ & $+16.25$ & $+16.25$ \\
    Person   & 61.25 &  86.25 &  98.12 &  99.38 & $+25.00$ & $+36.87$ & $+38.13$ \\
    \bottomrule
  \end{tabular}
\end{table}

\begin{table}[H]
  \caption{Two most wrong - Similar pattern}
  \label{tab:t2-confusions}
  \centering
  \footnotesize
  \begin{tabular}{llll}
    \toprule
    True class & Source-only & DANN & CDAN \\
    \midrule
    Dog      & elephant 333, giraffe 167 & person 397, house 18 & person 139, elephant 62 \\
    Horse    & giraffe 293, elephant 111 & dog 204, person 166  & dog 134, person 33 \\
    Elephant & person 21, giraffe 13     & person 400, dog 49   & person 65, dog 58 \\
    Giraffe  & elephant 99, horse 34     & person 164, dog 142  & dog 175, horse 56 \\
    \bottomrule
  \end{tabular}
  \\[2pt]
  \footnotesize Two most frequent wrong predictions per class, with counts.
\end{table}

Further data shows us that four classes show interesting behavior. Dog (+20.99 to +43.14) and Person (+25.00 to +38.13) improve under every method, while
Elephant and Giraffe degrade under every method. It almost seems like the model is trading the improvement in two for the others, no one improved all four.

\subsubsection{Controlled Design Study}

\begin{table}[H]
  \caption{Control study - too high a $lambda$ collapses the model}
  \label{tab:t2-study-compact}
  \centering
  \begin{tabular}{lccc}
    \toprule
    $\lambda_{\mathrm{MMD}}$ & Mean source F1 (\%) & Sketch acc.\ (\%) & Domain separability (\%) \\
    \midrule
    0 (source-only) & 94.85 & 61.59 & 100.00 \\
    0.1             & 94.73 & 62.31 &  92.17 \\
    1               & 94.93 & 64.44 &  81.87 \\
    10              &  5.07 &  4.07 &  60.03 \\
    \bottomrule
  \end{tabular}
\end{table}

A higher $\lambda$ means a higher penalty for difference between the source domain feature distribution, so increasing it should mean the model more aggressively tries
to match the feature distributions. And we see exactly that for the first three rows, however $\lambda = 10$ completely collapses the model. It is such a sever penalty
that all the distributions get so mixed up that the model can no longer tell any of them apart. F1 Mean on source alone is enough to disqualify 10 as optimial even without
consulting target domain labels.


\subsection{Open-Set Recognition (OSR)}

\subsubsection{Vanilla Closed-Set Baseline}

\subsection{Cross-Task Discussion}

\section{Conclusions}
\label{sec:conclusion}

\section{Appendix}
\label{sec:appendix}

\begin{figure}[t]
  \centering
  \includegraphics[width=\linewidth]{figures/tsne_resnet50.png}
  \caption{ResNet t-SNE: Representations overlap except in Cue Conflict}
  \label{fig:tsne-resnet}
\end{figure}

\begin{figure}[t]
  \centering
  \includegraphics[width=0.85\linewidth]{figures/tsne_vit_b_16.png}
  \caption{Vit-B t-SNE: Representations overlap except in Cue Conflict}
  \label{fig:tsne-vit}
\end{figure}

\begin{figure}[t]
  \centering
  \includegraphics[width=0.85\linewidth]{figures/tsne_clip.png}
  \caption{CLIP t-SNE: Representations overlap except in Cue Conflict and Patch Shuffle}
  \label{fig:tsne-clip}
\end{figure}

\section{References}
\label{sec:references}


\end{document}
